\label{seventh-bilet}
\begin{definition}
	\textit{Линейным нормированным пространством} над полем $\Kk$, где $\Kk = \R$ или $\Kk = \Cm$, называется линейное пространство $E$ над $\Kk$ с функцией $\|\cdot\| : E \to \R$, обладающей следующими свойствами:
	\begin{enumerate}
		\item $\forall x \in E: \|x\| \geqslant 0$, причем $\|x\| = 0 \lra x = 0$
		\item $\forall x \in E: \forall \alpha \in \Kk: \|\alpha x\| = |\alpha|\|x\|$
		\item $\forall x, y \in E: \|x + y\| \leqslant \|x\| + \|y\|$ (\textit{неравенство треугольника})
	\end{enumerate}

	Функция $\|\cdot\|$ называется \textit{нормой} на пространстве $E$.
\end{definition}

\begin{example}
    Легко видеть, что любое линейное нормированное пространство $E$ является метрическим с метрикой, заданной для произвольных $x, y \in E$ как $\rho(x, y) := \|x - y\|$. Приведенные ранее в качестве примеров метрические пространства являются линейными нормированными, и метрика в них имеет вид нормы разности элементов.
\end{example}

\begin{note}
	Норма непрерывна как функция $E \to \R$. Если для последовательности $\{x_n\} \subset E$ выполнено $x_n \xrightarrow[E]{} x$, то $\rho(x_n, x) = \|x_n - x\| \to 0$. Дважды воспользуемся неравенством треугольника:
	\begin{itemize}
		\item $\|x_n\| \leqslant \|x_n - x\| + \|x\|$
		\item $\|x\| \leqslant \|x_n - x\| + \|x_n\|$
	\end{itemize}

	Таким образом, $\left|\|x_n\| - \|x\|\right| \to 0$, то есть $\|x_n\| \to \|x\|$, что и требовалось.
\end{note}

\begin{definition}
	Пусть $E$ "--- линейное нормированное пространство. Множество $L \subset E$ называется:
	\begin{itemize}
		\item \textit{Линейным многообразием} в $E$, если $L$ замкнуто относительно сложения и умножения на скаляры из $\mathbb{K}$
  		\item \textit{Подпространством} в $E$, если $L$ является линейным многообразием в $E$ и при этом замкнуто 
	\end{itemize}
\end{definition}

\begin{definition}
	Пусть $E$ "--- линейное нормированное пространство, $M \subset E$. \textit{Линейной оболочкой} множества $M$ называется множество следующего вида:
	\[[M] := \left\{\sum_{k = 1}^n \alpha_im_i : \alpha_1, \dotsc \alpha_n \in \Kk, m_1, \dotsc, m_n \in M\right\}\]
\end{definition}

\begin{note}
	Линейная оболочка множества всегда является линейным многообразием в $E$, но не всегда является подпространством.
\end{note}

\begin{example}
    Пусть $E = C[0,1]$ и $M$ -- многочлены в $E$. Тогда $M = [M]$ является многообразием, но не является подпространством, так как линейной комбинацией многочленов можно представить только многочлен, но не все непрерывные функции.
    
    По теореме Вейерштрасса о приближении непрерывной функции многочленами, получается, что $[M]$ не содержит все свои точки прикосновения, а значит не замкнуто и не является подпространством.
\end{example}

\begin{definition}
	Пусть $E$ "--- линейное пространство. Нормы $\|\cdot\|$ и $\|\cdot\|'$ на $E$ называются \textit{эквивалентными}, если существуют числа $C_1, C_2 > 0$ такие, что для любого $x \in E$ выполнены неравенства:
	\[C_1\|x\|' < \|x\| < C_2\|x\|'\]
\end{definition}

\begin{note}
	Если нормы $\|\cdot\|$ и $\|\cdot\|'$ эквивалентны, то последовательность $\{x_n\} \subset E$ сходится или расходится относительно обеих норм одновременно.
\end{note}

\begin{definition}
	Пусть $E$ "--- линейное пространство. \textit{Размерностью} пространства $E$ называется наибольший размер линейно независимой системы в $E$. Обозначение "--- $\dim E$.
\end{definition}

\begin{proposition}[лемма о <<почти перпендикуляре>>]\label{almostperp}
	Пусть $E$ "--- линейное нормированное пространство, $L \subsetneq E$ "--- подпространство. Тогда для любого $\varepsilon > 0$ существует $y \in E$ такой, что $\|y\| = 1$ и $\rho(y, L) = \inf\limits_{z \in L}\|y - z\| \geqslant 1 - \varepsilon$.
\end{proposition}

\begin{proof}
		Зафиксируем $y_0 \in E \bs L$ положим $d := \rho(y_0, L) > 0$. Выберем вектор $z_0 \in L$ такой, что $d \leqslant \|y_0 - z_0\| \leqslant d(1 + \varepsilon)$ и покажем, что подходит вектор $y := \frac{y_0 - z_0}{\|y_0 - z_0\|}:= \alpha (y_0-z_0)$. Действительно, $\|y\| = 1$, и для любого $z \in L$ выполнены неравенства:
		\[\|y - z\| = \big\|\alpha(y_0 - z_0) - z\big\|=|\alpha| \big\| y_0 - \underbrace{\left(z_0 + \frac1{\alpha} z\right)}_{\in L} \big\| \geqslant |\alpha|d \geqslant \frac{\alpha}{\alpha(1+\varepsilon)} \geqslant 1 - \varepsilon\]

		Значит $\inf\limits_{z \in L}\|y - z\| \geqslant 1 - \varepsilon$. Получено требуемое.
\end{proof}

\begin{theorem}[Рисса]\label{thm4.1}
	Пусть $E$ "--- линейное нормированное пространство. Тогда единичная сфера $S(0, 1)$ компактна в $E$ $\lra$ $\dim E < +\infty$.
\end{theorem}

\begin{proof}~
	\begin{itemize}
		\item[$\la$] По \hyperref[normequiv]{эквивалентности норм в конечномерных пространствах}, все нормы на конечномерном линейном пространстве эквивалентны, а относительно евклидовой нормы сфера $S(0, 1)$ компактна.
  
  		\item[$\ra$] От противного. Зафиксируем $\varepsilon > 0$ и построим последовательность $\{x_n\} \subset S(0, 1)$ с попарными расстояниями не меньше $1 - \varepsilon$, из чего будет следовать, что сфера $S(0, 1)$ не вполне ограниченна и потому не компактна:
		\begin{itemize}
			\item[$\bullet$] Выберем $x_1 \in S(0, 1)$ произвольным образом
			\item[$\bullet$] По \hyperref[almostperp]{лемме о почти перпендикуляре} выберем $x_2 \in S(0, 1) \bs [x_1]$ такое, что $\rho(x_2, [x_1]) \geqslant 1- \varepsilon$, причем утверждение применимо так как линйеная оболочка конечномерна и по~\ref{closeness-subspace} она замкнута, а значит действительно подпространство.
			\item[$\bullet$] По \hyperref[almostperp]{лемме о почти перпендикуляре} выберем $x_3 \in S(0, 1) \bs [x_1, x_2]$ такое, что $\rho(x_3, [x_1, x_2]) \geqslant 1- \varepsilon$
		\end{itemize}

		Поскольку $\dim E = + \infty$, то процесс не закончится, и будет получена искомая последовательность $\{x_n\}$.\qedhere
	\end{itemize}
\end{proof}